\bookchapter{第六十二章：界面开发——SwiftUI 入门}

\chapterintro{界面是用户看到的一切。SwiftUI 是苹果 2019 年推出的\textbf{声明式}界面框架——你只需要“声明”界面长什么样，系统自动负责渲染、适配深色模式、适配各种屏幕。到 2026 年，SwiftUI 已是新项目的绝对主流。这一章，咱们学会“画画”。}

\booksection{一、SwiftUI 的思想：声明式，不是命令式}

老式界面开发是“命令式”：先告诉系统“创建按钮 放在(20,30) 设置颜色”，一步一步指挥。SwiftUI 是“声明式”：\textbf{你直接描述“界面上有一个红色按钮写着'点我'”}，怎么画、什么时候重画，系统自己操心。

\begin{lstlisting}[language=Swift]
import SwiftUI

struct ContentView: View {
    var body: some View {      // body：描述"界面长什么样"
        VStack {               // 竖着排列
            Text("你好，世界！")
                .font(.largeTitle)      // 修饰符：大标题
            Button("点我") {
                print("被点击了")
            }
            .buttonStyle(.borderedProminent)  // 蓝色实心按钮
        }
        .padding()
    }
}
\end{lstlisting}

关键点：\textbf{视图 = 遵守 View 协议的结构体}，\texttt{body}属性描述界面，\texttt{.\allowbreak 修饰符()}链式“上妆”。这就是 SwiftUI 的全部基本盘。

\booksection{二、基础视图：文字、按钮、输入框、图片}

\begin{center}
\begin{tabularx}{\linewidth}{|X|X|X|}
\hline
\textbf{视图} & \textbf{干什么} & \textbf{示例} \tabularnewline
\hline
Text & 显示文字 & \texttt{Text("你好")} \tabularnewline
Button & 按钮 & \texttt{Button("点我") \{ 动作 \}} \tabularnewline
TextField & 输入框 & \texttt{TextField("提示", text: \$文本)} \tabularnewline
Image & 显示图片 & \texttt{Image(systemName: "star.\allowbreak fill")} \tabularnewline
\hline
\end{tabularx}
\end{center}

\begin{center}\small\color{gray}SwiftUI 四大基础视图\end{center}

注意\texttt{Image(systemName:)}——苹果内置了\textbf{SF Symbols}图标库（几千个免费图标），“star.fill” 这种名字就能直接显示一个实心星星。这是 iOS 开发的隐藏福利：\textbf{做界面几乎不用找图标素材}。

\booksection{三、布局三件套：VStack、HStack、ZStack}

\begin{itemize}
  \item \textbf{VStack}：竖着排（像排队打饭）；
  \item \textbf{HStack}：横着排（像并排坐）；
  \item \textbf{ZStack}：叠着放（像扑克牌摞一起，适合做角标、遮罩）。
\end{itemize}

\begin{lstlisting}[language=Swift]
VStack(alignment: .leading, spacing: 12) {
    HStack {
        Text("账单")
        Spacer()                 // 弹性空间：把右边推到尽头
        Text("-25.5 元").foregroundColor(.red)
    }
    Divider()                    // 分割线
    Text("备注：午饭")
}
.padding(16)                     // 四周留白
.background(Color(.systemGray6)) // 浅灰底
.cornerRadius(12)                // 圆角

// 修饰符顺序有意义：先 padding 再 background，背景不含边距
\end{lstlisting}

\begin{tcbwarn}{注意：修饰符顺序有讲究}
\texttt{padding(16).\allowbreak background(灰)}和\texttt{background(灰).\allowbreak padding(16)}效果不同——前者“先留白再涂色”（背景不含边距），后者“先涂色再留白”（背景包含边距）。新手最常见的“奇怪白边”都源于此。不确定时调换顺序，预览立刻告诉你答案。

\end{tcbwarn}

\booksection{四、状态与绑定：让界面“活”起来}

静态界面没意思。SwiftUI 的规则：\textbf{界面永远忠实反映状态，状态变了，界面自动重画}。

\begin{lstlisting}[language=Swift]
struct CounterView: View {
    @State private var count = 0   // @State：这个值变了要重画

    var body: some View {
        VStack {
            Text("你点了 \(count) 次")
            Button("点我 +1") {
                count += 1          // 改状态，界面自动刷新
            }
        }
    }
}
\end{lstlisting}

输入框要用“双向绑定”——\texttt{\$}符号把“值和修改权”一起传进去：

\begin{lstlisting}[language=Swift]
struct InputView: View {
    @State private var text = ""

    var body: some View {
        TextField("输入点什么", text: $text)  // $ 双向绑定
            .textFieldStyle(.roundedBorder)
        Text("你输入了：\(text)")   // 输入时实时更新
    }
}
\end{lstlisting}

用户操作

状态改变（count += 1）

SwiftUI 感知到变化

自动重画界面（body 重新计算）

\begin{center}\small\color{gray}图 3：SwiftUI 的“状态 界面”单向数据流\end{center}

\booksection{五、预览、主题与动画：SwiftUI 的“外挂”}

\begin{itemize}
  \item \textbf{\#Preview 宏}：文件底部加\texttt{\#Preview \{ ContentView() \}}，右侧画布实时渲染——改代码所见即所得；
  \item \textbf{深色模式免费}：用系统颜色（\texttt{Color(.\allowbreak systemGray6)}、\texttt{.\allowbreak primary}）而不是写死颜色，深色模式自动适配，一行不用多写；
  \item \textbf{动画一行搞定}：\texttt{withAnimation \{ count += 1 \}}或给视图加\texttt{.\allowbreak animation(.\allowbreak default, value: count)}——过渡丝滑。
\end{itemize}

\begin{lstlisting}[language=Swift]
struct AnimatedHeartView: View {
    @State private var liked = false

    var body: some View {
        Button {
            withAnimation(.spring(duration: 0.4)) {
                liked.toggle()
            }
        } label: {
            Image(systemName: liked ? "heart.fill" : "heart")
                .font(.system(size: 48))
                .foregroundColor(liked ? .red : .gray)
                .scaleEffect(liked ? 1.3 : 1.0)
        }
    }
}
\end{lstlisting}

\begin{tcbinfo}{小贴士：先功能后动画}
新手常犯的毛病是“界面还没做对，就开始纠结动画”。先让功能跑通，最后再统一加动效——动画是锦上添花，功能是雪中送炭。这也是咱们一贯的“简单至上”原则。

\end{tcbinfo}

\begin{tcbpractice}{来动手：做一个"心情按钮"}
\booksection{六、常用修饰符速查：一页在手，界面不愁}

\begin{center}
\begin{tabularx}{\linewidth}{|X|X|}
\hline
\textbf{想做什么} & \textbf{修饰符} \tabularnewline
\hline
留白 & \texttt{.\allowbreak padding(16)} \tabularnewline
圆角 & \texttt{.\allowbreak cornerRadius(12)} \tabularnewline
背景色 & \texttt{.\allowbreak background(Color(.\allowbreak systemGray6))} \tabularnewline
字体与字号 & \texttt{.\allowbreak font(.\allowbreak title.\allowbreak bold())}/\texttt{.\allowbreak font(.\allowbreak system(size: 20))} \tabularnewline
前景色（文字/图标） & \texttt{.\allowbreak foregroundColor(.\allowbreak red)} \tabularnewline
撑满宽度 & \texttt{.\allowbreak frame(maxWidth: .\allowbreak infinity)} \tabularnewline
固定尺寸 & \texttt{.\allowbreak frame(width:100, height: 44)} \tabularnewline
对齐方式 & \texttt{VStack(alignment: .\allowbreak leading)} \tabularnewline
透明度 & \texttt{.\allowbreak opacity(0.\allowbreak 6)} \tabularnewline
阴影 & \texttt{.\allowbreak shadow(radius: 5)} \tabularnewline
点击手势 & \texttt{.\allowbreak onTapGesture \{ .\allowbreak .\allowbreak .\allowbreak  \}} \tabularnewline
键盘类型 & \texttt{.\allowbreak keyboardType(.\allowbreak decimalPad)} \tabularnewline
\hline
\end{tabularx}
\end{center}

\begin{center}\small\color{gray}SwiftUI 常用修饰符速查\end{center}

记不住没关系——这页就是用来查的。用得最多的其实就前四个：padding、圆角、背景、字体。

把模板项目改成一个心情切换器：一个按钮，点击后在 开心 难过 生气 之间循环切换，文字和背景色跟着变。提示：用数组存三组（文字+颜色），\texttt{@State}存下标，点击时\texttt{(index + 1) \% 3}。做完发给朋友玩，收获第一波“哇你会做 App 了”。

\end{tcbpractice}

\begin{tcbsummary}{本章小结}
\begin{itemize}
  \item SwiftUI 是声明式：描述"长什么样"，系统负责画和重画
  \item 视图 = 遵守 View 协议的结构体；body 描述界面；修饰符链式"上妆"
  \item 四大基础视图：Text / Button / TextField / Image；SF Symbols 图标库免费管饱
  \item 布局三件套：VStack 竖排、HStack 横排、ZStack 叠放；Spacer 是弹性空间
  \item @State 让界面活起来；\$ 双向绑定喂给输入框
  \item \#Preview 所见即所得；系统颜色自动适配深色模式；withAnimation 一行动效
\end{itemize}

\end{tcbsummary}
